\subsection{Лекция 18 (07.03.26) Coupled решатель на основе интерполяции Rhie--Chow}
В тесте \cvar{[cavity-fvm-coupled]} из файла \ename{cavity_fvm_coupled_test.cpp}
реализована Coupled схема решения задачи о течении в каверне.
Стационарное решение получено методом установления.
Неявное конвективное слагаемое собрано с применением MC ограничителя (по методике \secref{sec:grad_inteprolate_up}).
У расчётного алгоритма есть единственный параметр -- шаг по времени $\tau$.
Необходимо:
\begin{enumerate}
\item  Нарисовать раздельно графики сходимости для 
       $$\max\left|\dfr{u}{t}\right|,\quad
         \max\left|\dfr{v}{t}\right|,\quad
         \max\left|\dfr{p}{t}\right|$$
\item  Исследовать схему на оптимальный $\tau$ с точки зрения скорости сходимости с изменением шага по пространству.
       Провести аналогичный анализ для других чисел Рейнольдса (до 1000). Проверить полученные оптимумы на неструктурированной pebi сетке.
\item  Сравнить решения полученные с применением MC ограничителя с решением без ограничителя при разных числах Рейнольдса.
\item  Провести профилировку и рефакторинг программы, убрав повторяющийся код и выделив процедуры, которые можно выполнять один раз на этапе инициализации.
\item  Модифицировать алгоритм, отнеся конвективные слагаемые в уравнениях моментов полностью на прошлый временной слой.
       Проверить, как при этом изменится скорости расчёта и скорости сходимости. Нарисовать сравнительные графики.
\end{enumerate}

\paragraph{Рекомендации к программированию}
\begin{itemize}
\item п. 3
Для отключения ограничителей достаточно, чтобы метод
\cvar{tvd_upwind_weight} всегда возрвращал $0.5$.
\item п. 4
Для профилирования можно воспользоваться методами, описанными в \secref{sec:dbg-tictoc}.
\item п. 4
При рефакторинге обратить внимание на
\begin{itemize}
\item Повторяющийся код в методах \cvar{u_equation}, \cvar{v_equation},
\item Повторяющийся алгоритм при вычислении поправки Rhie-Chow в методах \cvar{continuity_equation}, \cvar{compute_un_rhie_chow},
\item Некоторые части блочной матрицы не зависят от временного слоя и могли бы быть собраны один раз
\end{itemize}
При рефакторинге необходимо следить за неизменностью резульатов тестовых проверок.
\item п. 5
Обратить внимание, что использование явной конвекции делает всю матрицу левой части неизменной.
Тогда весь решатель можно собрать и инициализировать один раз.
\end{itemize}
